\Needspace{10\baselineskip}
\Article{10}
Later, the celebrated \textsc{de Foncenex} too, having noticed a defect in \textsc{Euler's} first proof (above, Article 8, objection 4) that he could not remove, tried another path and presented it on p. 120 of the paper cited\HistoricalNote{7}. It consists in the following.

Let an equation $Z=0$ be proposed, where $Z$ denotes a function of degree $m$ in the unknown $z$. If $m$ is an odd number, it is already known that this equation has a real root; if $m$ is even, however, the celebrated Foncenex tries in the following way to prove that the equation has at least one root of the form $p+q\BareRoot{-1}$. Let $m=2^n i$, where $i$ denotes an odd number, and suppose that $zz+uz+M$ is a divisor of the function $Z$. Then the individual values of $u$ will be the sums of pairs of roots of the equation $Z=0$, with their sign changed; hence $u$ will have $\frac{m\mathbin{.}m-1}{1\mathbin{.}2}=m'$ values. If $u$ is supposed to be determined by an equation $U=0$ (where $U$ denotes a polynomial in $u$ and the known coefficients in $Z$), this equation will have degree $m'$. One readily sees that $m'$ will be a number of the form $2^{n-1}i'$, where $i'$ denotes an odd number. Now, unless $m'$ is odd, suppose again that $uu+u'u+M'$ is a divisor of $U$; similar reasoning will show that $u'$ is determined by an equation $U'=0$, where $U'$ is a function of degree $\frac{m'\mathbin{.}m'-1}{1\mathbin{.}2}$ in $u'$. Putting $\frac{m'\mathbin{.}m'-1}{1\mathbin{.}2}=m''$, we have that $m''$ is a number of the form $2^{n-2}i''$, where $i''$ denotes an odd number. Now, unless $m''$ is odd, let $u'u'+u''u'+M''$ be a divisor of the function $U'$; then $u''$ will be determined by an equation $U''=0$, which, if it is supposed to have degree $m'''$, gives $m'''$ as a number of the form $2^{n-3}i'''$. It is evident that in the sequence of equations $U=0$, $U'=0$, $U''=0$, etc., the one in position $n$ will have odd degree and therefore a real root. For brevity we shall put $n=3$, so that the equation $U''=0$ has a real root $u''$, for one readily sees that the same reasoning holds for any other value of $n$. The celebrated de Foncenex then asserts that the coefficient $M''$ can be determined rationally from $u''$ and the coefficients in $U'$ (which are readily understood to be polynomials in the coefficients in $Z$), or from $u''$ and the coefficients in
\SourceLeaf{024}
$Z$, and is therefore real. It follows that the roots of the equation $u'u'+u''u'+M''=0$ will be contained in the form $p+q\BareRoot{-1}$; but they will evidently satisfy the equation $U'=0$ as well. There will therefore be some value of $u'$ contained in the form $p+q\BareRoot{-1}$. Now the coefficient $M'$ can, in the same way as before, be determined rationally from $u'$ and the coefficients in $Z$, and will therefore also be contained in the form $p+q\BareRoot{-1}$. Hence the roots of the equation $uu+u'u+M'$ will be contained in the same form, and will at the same time satisfy the equation $U=0$; that is, this equation will have a root contained in the form $p+q\BareRoot{-1}$. Finally, it follows from this in a similar way that $M$ too is contained in the same form, as is a root of the equation $zz+uz+M=0$, which will evidently also satisfy the proposed equation $Z=0$. Thus any equation will have at least one root of the form $p+q\BareRoot{-1}$.

\Article{11}
Objections 1, 2 and 3 that I made against \textsc{Euler's} first proof (Article 8) have the same force against this method, with this difference, however: the second objection, to which \textsc{Euler's} proof was open only in some special cases, must affect the present one in every case. Namely, one can prove a priori that, even if a formula is given that expresses the coefficient $M'$ rationally in terms of $u'$ and the coefficients in $Z$, this formula must necessarily become indeterminate for several values of $u'$; similarly, the formula that expresses the coefficient $M''$ in terms of $u''$ becomes indeterminate for some values of $u''$, etc. This is seen most clearly if we take an equation of the fourth degree as an example. Let us therefore put $m=4$, and let the roots of the equation $Z=0$ be $\alpha$, $\beta$, $\gamma$, $\delta$. Then it is evident that the equation $U=0$ will be of the sixth degree and its roots $-(\alpha+\beta)$, $-(\alpha+\gamma)$, $-(\alpha+\delta)$, $-(\beta+\gamma)$, $-(\beta+\delta)$, $-(\gamma+\delta)$. The equation $U'=0$, however, will be of the fifteenth degree, and the values of $u'$ will be
\[
\begin{gathered}
2\alpha+\beta+\gamma,\quad 2\alpha+\beta+\delta,\quad 2\alpha+\gamma+\delta,\\
2\beta+\alpha+\gamma,\quad 2\beta+\alpha+\delta,\quad 2\beta+\gamma+\delta,\\
2\gamma+\alpha+\beta,\quad 2\gamma+\alpha+\delta,\quad 2\gamma+\beta+\delta,\\
2\delta+\alpha+\beta,\quad 2\delta+\alpha+\gamma,\quad 2\delta+\beta+\gamma,\\
\alpha+\beta+\gamma+\delta,\quad\alpha+\beta+\gamma+\delta,\quad\alpha+\beta+\gamma+\delta
\end{gathered}
\]
Now we must stop at this equation, since its degree is odd, and it will in fact have the real root $\alpha+\beta+\gamma+\delta$ (which equals the first coefficient in $Z$ with its sign changed and will therefore be not only real but also rational if the coeffi\SourceLeaf{025}cients in $Z$ are rational). But one can readily see that, if a formula is given that expresses the value of $M'$ rationally in terms of the corresponding value of $u'$, it necessarily becomes indeterminate for $u'=\alpha+\beta+\gamma+\delta$. For this value will be a root of the equation $U'=0$ \textit{three times}, and three values of $M'$ will correspond to it, namely $(\alpha+\beta)(\gamma+\delta)$, $(\alpha+\gamma)(\beta+\delta)$ and $(\alpha+\delta)(\beta+\gamma)$, which can all be irrational. Evidently a rational formula could in this case produce neither an irrational value of $M'$ nor three different values. From this example one can sufficiently gather that the celebrated \textsc{de Foncenex's} method is by no means satisfactory; if it is to be made complete in every respect, the theory of elimination must be investigated much more deeply.

\Needspace{16\baselineskip}
\Article{12}
Finally, the distinguished \textsc{La Grange} treated our theorem in the paper \textit{Sur la forme des racines imaginaires des équations, Nouv. Mém. de l'Acad. de Berlin} 1772, p. 222 ff. This great mathematician took particular pains to supply the defects in \textsc{Euler's} first proof, and in fact investigated especially those matters that constitute the second and fourth objections above (Article 8) so deeply that nothing further remains to be desired, except perhaps that some doubts seem to remain in the earlier investigation concerning the theory of elimination, on which this whole investigation rests. --- Yet he did not touch on the third objection at all; indeed, the whole investigation is built on the supposition that any equation of degree $m$ actually has $m$ roots.

Thus, after carefully weighing what has been set out so far, I hope that those familiar with the subject will welcome a new proof of this most important theorem drawn from entirely different principles. I now proceed to give it.

\Article{13}
\textsc{Lemma.} \textit{If $m$ denotes any positive integer, the function $\sin\varphi\mathbin{.}x^m-\sin m\varphi\mathbin{.}r^{m-1}x+\sin(m-1)\varphi\mathbin{.}r^m$ will be divisible by $xx-2\cos\varphi\mathbin{.}rx+rr$.}

\textit{Proof.} For $m=1$ this function becomes $=0$ and is therefore divisible by any factor; for $m=2$ the quotient becomes $\sin\varphi$, and for any greater value the quotient will be $\sin\varphi\mathbin{.}x^{m-2}+\sin2\varphi\mathbin{.}rx^{m-3}+\sin3\varphi\mathbin{.}rrx^{m-4}+\text{etc.}+\sin(m-1)\varphi\mathbin{.}r^{m-2}$. For one readily verifies that, when this function is multiplied by $xx-2\cos\varphi\mathbin{.}rx+rr$, the product equals the proposed function.
